% Section 7: Conclusion

\section{Conclusion}
\label{sec:implications}

The FS~AI~RMF gives financial institutions a substantial control vocabulary
\parencite{treasuryFSAIRMFRelease2026}. It does not, by itself, tell a community
bank how those controls cohere or tell an examiner what a completed governance
artifact warrants. Empty-chair representation offers a proposed organizing
lens: identify the party absent from an AI-supported decision, the interest a
control is meant to represent, and the evidence by which that representation
can later be tested. The worked examples show how to apply the lens to selected
control objectives; framework-wide validation remains future work.

The paper's central distinction is between an artifact and the inference drawn
from it. Decision-time records, feature attributions, local approximations,
internal-state observations, generated rationales, human deliberative records,
and institutional approvals can all be useful. They expose different
evidentiary objects. An explanation reviewed by a human establishes that the
explanation was reviewed; stronger conclusions depend on what the method
targets, what the reviewer could test, and how the evidence is bound to the
claim. This distinction narrows neither AI nor explanation to a single approved
form. It requires institutions to state the proposition each artifact is being
used to support.

Silence-manufacture names a recurring inference gap: relevant evidence is
unavailable, a different artifact is offered, and the absence of contradictory
evidence supports a conclusion the artifact does not independently establish.
The pattern is a diagnostic, not a presumption of bad faith and not a claim that
every governance artifact fails. It directs attention to the relation among the
proposition, the artifact, the claimed binding, and the party who cannot test
that binding.

The three diagnostic techniques turn that distinction into a shared inquiry.
Empty-chair enumeration identifies whose interest depends on an
artifact. Pages missing from the ledger makes declared and ambient evidentiary
absences visible. The produced-silence question asks what access a design
removed and what is inferred from the resulting gap. Provisional evidentiary
capabilities then identify properties that may help participants answer those
questions. Vendors and architects are generally positioned to implement such
capabilities; banks can identify gaps, use compensating controls, and express
procurement needs; examiners can test claims within current authority; and
regulators can decide whether recurring gaps warrant coordinated expectations.
Their implementation, sufficiency, and tractability remain to be tested.

For a participant reading or contesting an AI governance record, the frame
reduces to five questions:

\begin{enumerate}
  \item What evidentiary object does this artifact expose?
  \item What conclusion is the institution asking the examiner to draw from it?
  \item What binds the artifact to that conclusion?
  \item Whose interest bears the cost if the binding is absent?
  \item What additional evidence would close or honestly disclose the gap?
\end{enumerate}

Those questions do not decide every governance dispute. They do make the point
of disagreement inspectable. A useful AI governance architecture should let an
absent party's interest appear not only in policy, but in evidence capable of
bearing the conclusion placed upon it.
